% Generated draft: requires mathematical and source review.
\documentclass[11pt]{article}
\usepackage[margin=25mm]{geometry}
\usepackage{fontspec}
\setmainfont{Latin Modern Roman}
\usepackage{amsmath,amssymb,mathtools,unicode-math}
\setmathfont{Latin Modern Math}
\usepackage{enumitem,needspace,xurl,hyperref}
\setlength{\parindent}{0pt}
\setlength{\parskip}{5pt}
\setlength{\emergencystretch}{3em}
\newcommand{\N}{\mathbb{N}}
\newcommand{\Z}{\mathbb{Z}}
\newcommand{\Q}{\mathbb{Q}}
\newcommand{\R}{\mathbb{R}}
\newcommand{\C}{\mathbb{C}}
\newcommand{\F}{\mathbb{F}}
\newcommand{\A}{\mathbb{A}}
\newcommand{\PP}{\mathbb{P}}
\newcommand{\E}{\mathbb{E}}
\newcommand{\eps}{\varepsilon}
\newcommand{\dd}{\,\mathrm{d}}
\begin{document}
\section*{UnsolvedMath ID: 8800084}
\textbf{Source problem code:} catalogue entry
\subsection{Definitions and mathematical statement}
Algebraically, a bilinear pairing between finite abelian groups can be written
\[
e:G_1\times G_2\to G_T,\qquad
 e(aP,bQ)=e(P,Q)^{ab},
\]
with nondegeneracy specified for the intended groups. A trapdoor construction additionally has public parameters \(\mathrm{pp}\) and secret information \(\tau\), generated together, such that some precisely identified operation is efficient with \(\tau\) and is assumed or proved hard from the public data alone.

Rivest asks for improved practical trapdoor pairings. The source does not state what the trapdoor enables: evaluating an otherwise hidden pairing, inverting a publicly evaluable pairing, recognizing hidden subgroup structure, or another task are different interfaces and security problems. It therefore does not authorize choosing one of them as the definition. A generic setup notation
\[
\operatorname{Gen}(1^\lambda)\longmapsto(\mathrm{pp},\tau)
\]
only identifies where the public/secret data would enter; correctness and the secret-enabled algorithm remain to be specified.

A concrete answer must fix those operations and their security experiment, group encodings and order structure, parameter distribution, and performance criteria. Improved needs a baseline construction, while practical needs target key sizes, computation/memory costs, and security parameters. No such comparison data or formal hardness assumption are included. The usual existence of public bilinear pairings alone does not provide the intended trapdoor feature.

The supplied research summary reports no verified practical improvement and describes specialized proposals as unverified/impractical or unpublished. This is a search-limited assessment, not a nonexistence theorem or an assertion that all conceivable trapdoor interfaces are equivalent.
\subsection{Short English statement}
Construct a practically improved pairing with a precisely specified secret trapdoor advantage and security proof. The secret-enabled operation, baseline construction, and performance/security targets are absent from the record.
\subsection{Sources}
\begin{itemize}
\item Local UnsolvedMath registry, permanent ID \(8800084\), exact source code \texttt{catalogue entry}, ``Better trapdoor pairings.'' Proposer Ronald Rivest.
\item IPAM \emph{Workshop on Number Theory and Cryptography --- Open Problems} (2006), attachment 6243, legacy PPT logical slide 17; \url{https://www.ipam.ucla.edu/programs/workshops/workshop-i-number-theory-and-cryptography-open-problems/?tab=schedule}. Recorded access 2026-07-29; pdf-and-legacy-ppt visual-audit extraction.
\item Supplied OPEN-TRIAGE research summary mentions specialized proposals but gives no construction name, exact trapdoor definition, or publication citation.
\end{itemize}
\end{document}
